\scanpage{182}
% archive sheet number: 181
sur $ {\Gamma}^+$, c'est la fonction ``dimension d'un
module''. Dans le cas \uline{où} on a un $e$,
on a $n(e)=1$, et si on a un $\unsure{\nu|v|\vee}$, on a
$n(\check i)=n(i)$.

$\unsure{\mathcal K|\mathscr K|\mathcal R}$ est bien définie si, en plus de
$I$ et de la fonction $n(i)$, on se donne un
représentant $M_i$ de chaque $i$ avec comme
module sous-jacent $k^{n(i)}$, et pour $i,j$ donnés,
un sous-groupe $\operatorname{Hom}_{j,\deleted{\unsure{j|k}}i}(k^{n(i)},k^{n(j)})$,
de façon que $\operatorname{Hom}_{\unsure{kj|kj\,j}}\times\operatorname{Hom}_{\unsure{ji|ij}}$ soit appliqué
par composition dans $\operatorname{Hom}_{ki}$, etc\ldots.

Supposons que dans $K$ on ait une
notion de produit tensoriel, alors
$\Gamma(K)$ est muni d'une loi d'anneau \inserted{associatif}
et commutatif, qui sur $\Gamma^+(K)$ \deleted{\unsure{p|y}} correspond
à la loi de produit tensoriel. $\underline n$ est une
fonction multiplicative sur $\Gamma^+$, et la
compatibilité avec $e$ \deleted{signifie} \inserted{implique} que $e$ est
une unité pour \deleted{\unsure{l'op}} la loi d'algèbre. La
compatibilité avec $\nu$ \deleted{signifie} \inserted{implique} que $\nu$ est
un automorphisme de l'algèbre $\Gamma$. Enfin
la compatibilité simultanée de
$e,\nu,\otimes$ implique que $i*\check i$ a toujours
\deleted{\unsure{p|pour}}pour composante non nulle suivant $e$.
$I$ est alors un ``\uline{hypergroupe discret}''
